\documentclass[11pt,a4paper]{article}
\usepackage[margin=25mm]{geometry}
\usepackage[T1]{fontenc}
\usepackage{lmodern}
\usepackage{amsmath,amssymb,amsthm,mathtools}
\usepackage{microtype}
\usepackage[hidelinks]{hyperref}
\newtheorem{theorem}{Theorem}
\newtheorem{corollary}[theorem]{Corollary}
\newtheorem{remark}[theorem]{Remark}
\newcommand{\N}{\mathbb N}
\title{An Extreme-Spike Barrier for the Collatz Clipping Tail\\
\large An internal consequence of odd-core compression}
\author{Research checkpoint --- not a proof of the asymptotic tail bound}
\date{17 November 2026}
\begin{document}
\maketitle
\begin{abstract}
We derive a pointwise, one-sided sufficient condition for the vanishing
of the clipping tail from the \emph{odd-core compression} estimate and the
sparsity of dangerous parity cylinders already established in the
consolidated manuscript and its companion note. A failure of
$T_X=o(X)$ forces exceptional odd target weights of size at least
$\exp(\kappa (\log X)^{\varepsilon_0})$ on an infinite sequence of scales.
This is a rigorous \emph{necessary condition for failure}, not a proof
that these spikes do not exist. The required pointwise estimate for the
true logarithmic predecessor potential remains open.
\end{abstract}

\section{Inputs and notation}
Fix the accelerated map $T_{3,c}$, odd $c\ge-1$, $3\nmid c$, and the
constants of the consolidated manuscript, in particular $\beta=1/200$
and $0<\varepsilon_0<1$. Put $X_1=X^{1-2\beta}$.
Let $h_X^\downarrow(a)=h(a)$ if the forward orbit of $a$ reaches
$[1,X_1]$, and $h_X^\downarrow(a)=0$ otherwise.  Set
\[
 g_X(a):=a h_X^\downarrow(a),\qquad
 \mu_X:=\mu_X(L_X),\qquad
 \mu_X\le C_\mu(\log X)^{-\varepsilon_0},
\]
where $L_X=(\log X)^{\varepsilon_0/2}$ and $C_\mu>0$ is fixed.
We do not reprove the first-passage inputs here.

Fix $\zeta_0=0.55$ and define
\[
 \zeta=\frac{\log2}{\log3},\qquad
 \Delta=\zeta_0\log(2\zeta_0)+(1-\zeta_0)\log(2-2\zeta_0)>0,
 \qquad t_X=\left\lceil\frac{3\beta\log X}{(\zeta-\zeta_0)\log3}\right\rceil.
\]
Numerically $\Delta\approx0.00500837$ and
$2^{t_X}\le 2X^{0.12}$ for all sufficiently large $X$.
For $b_0\le b\le t_X$, let $\mathcal C_{X,b}$ denote the union
of residue classes modulo $2^{t_X}$ whose parity word has at least
$\zeta_0 t$ odd entries in some prefix $b\le t\le t_X$.
As established by parity-word counting,
\begin{equation}\label{eq:sparse}
 \#(\mathcal C_{X,b}\cap[1,X])\le X\alpha_b+2^{t_X},
 \qquad \alpha_b:=\frac{e^{-\Delta b}}{1-e^{-\Delta}}.
\end{equation}

Our sole dynamical input beyond these definitions is the previously
proved odd-core bound (companion note, Theorem ``Odd-core upper bound
for the clipping tail''):
\begin{equation}\label{eq:oddinput}
 \frac{T_X}{X}\le (5+6b)\mu_X+
 \frac{36}{X}\sum_{\substack{X_1/4<a\le X\\a\text{ odd},\ a\in\mathcal C_{X,b}}}g_X(a)
 \quad(b_0\le b\le t_X).
\end{equation}
Thus the claims below are conditional on the first-passage and odd-core
propositions, not independent of their hypotheses.

\section{Pointwise extreme-spike theorem}
For each large $X$ define the \emph{short-valuation odd maximum}
\[
 G_X:=\max\Bigl(\{0\}\cup
 \bigl\{g_X(a):\ X_1/4<a\le X,\ a\text{ odd},\
 \nu_2(3a+c)\le\lfloor(1-\zeta_0)t_X+1\rfloor\bigr\}\Bigr).
\]
This maximum is finite and $0\le G_X\le X$. The valuation restriction
is justified by the proved deterministic initial-valuation lemma:
every odd $a\in\mathcal C_{X,b}$ satisfies that restriction for $b\ge2$.

\begin{theorem}[Quantitative extreme-spike bound]\label{thm:main}
Under the stated inputs, for all large $X$ and all
$\max(2,b_0)\le b\le t_X$,
\begin{equation}\label{eq:main}
 \boxed{\quad
 \frac{T_X}{X}\le
 (5+6b)C_\mu(\log X)^{-\varepsilon_0}
 +36G_X\left(\frac{e^{-\Delta b}}{1-e^{-\Delta}}
 +2X^{-0.88}\right).\quad}
\end{equation}
In particular:
\begin{enumerate}
\item If $\log(2+G_X)=o((\log X)^{\varepsilon_0})$ as
$X\to\infty$, then $T_X=o(X)$.
\item If for some fixed $A\ge0$ one has
$G_X\le(\log X)^A$ for all sufficiently large $X$, then
\begin{equation}\label{eq:polylog}
 \frac{T_X}{X}
 =O_{A,c}\bigl((\log\log X)(\log X)^{-\varepsilon_0}
 +(\log X)^{-2}+X^{-0.88}(\log X)^A\bigr).
\end{equation}
\item Conversely, if $\limsup_{X\to\infty}T_X/X>0$, then
there are constants $\kappa>0$ and $\delta>0$ and a sequence
$X_j\to\infty$ such that
\begin{equation}\label{eq:spike}
 \frac{T_{X_j}}{X_j}\ge\delta,
 \qquad
 G_{X_j}\ge\exp\bigl(\kappa(\log X_j)^{\varepsilon_0}\bigr).
\end{equation}
Thus failure of the clipping-tail conclusion forces
super-polylogarithmic odd-target spikes \emph{among stopped
orbits with short initial valuation}.
\end{enumerate}
\end{theorem}
\begin{proof}
By the initial-valuation lemma, every odd state counted in
\eqref{eq:oddinput} has weight at most $G_X$. The number of such
states is at most $X\alpha_b+2^{t_X}$, by \eqref{eq:sparse}.
Using $2^{t_X}\le2X^{0.12}$ proves \eqref{eq:main}.

For the first consequence put $L=(\log X)^{\varepsilon_0}$ and
$H_X=\log(2+G_X)=o(L)$. Choose
\[
 b=\max\left(\lceil b_0\rceil,2,
 \left\lceil\frac{H_X+\log\log X}{\Delta}\right\rceil\right).
\]
Because $0<\varepsilon_0<1$, we have $b=o(L)=o(\log X)$, so
$b\le t_X$ for large $X$. The first term in \eqref{eq:main}
is $o(1)$. The exponential term is at most a constant multiple
of $(\log X)^{-1}$, while
$G_X X^{-0.88}=\exp(o(\log X))X^{-0.88}=o(1)$.
This proves (1).

For (2) choose
$b=\max(\lceil b_0\rceil,2,
\lceil (A+2)\Delta^{-1}\log\log X\rceil)$.
This is below $t_X$ for all sufficiently large $X$.
Substitution in \eqref{eq:main} gives \eqref{eq:polylog}.

For (3), select $\delta>0$ and an infinite sequence of $X$
with $T_X/X\ge\delta$. Set $L=(\log X)^{\varepsilon_0}$ and
$b=\lfloor\delta L/(24C_\mu)\rfloor$. Eventually
$\max(2,b_0)\le b\le t_X$, and
$(5+6b)C_\mu/L\le\delta/3$.
Thus \eqref{eq:main} yields
\[
 \frac{2\delta}{3}\le
 36G_X\left(\frac{e^{-\Delta b}}{1-e^{-\Delta}}
 +2X^{-0.88}\right).
\]
Since $b=O((\log X)^{\varepsilon_0})=o(\log X)$,
we have $X^{-0.88}=o(e^{-\Delta b})$. Therefore
$G_X\ge c_\delta e^{\Delta b}$ for a positive constant
$c_\delta$, proving \eqref{eq:spike} with any
$0<\kappa<\Delta\delta/(24C_\mu)$ on the sequence.
\end{proof}

\section{Interpretation and limits}
\begin{remark}[A genuine one-way criterion, not a completed proof]
The condition in Theorem~\ref{thm:main}(1) is a \emph{sufficient}
pointwise upper bound.  Its converse is not claimed: large isolated
values of $G_X$ need not force a large clipping tail.  It is
logically different from the previously criticised concentration
criteria that are equivalent to the tail statement.  It also differs
from a uniform $L^p$ condition: an abstract profile that takes the
value $(\log X)^A$ on a positive proportion of the odd targets
satisfies the pointwise criterion but has divergent normalized $L^p$
means for every $p>1$ when $A>0$. Such an abstract profile need
not satisfy the Collatz harmonic equations.
\end{remark}
\begin{remark}[The hard missing input]
No theorem here establishes
$\log(2+G_X)=o((\log X)^{\varepsilon_0})$ for the true
logarithmic predecessor potential. The trivial inequality
$h_X^\downarrow(a)\le1$ gives only $G_X\le X$, far weaker
than required. Neither parity-word counting (Terras) nor the
available first-passage total-variation estimates (Tao/GGM)
supplies this pointwise estimate. The constant
$\Delta\approx0.005$ also makes the explicit asymptotic bounds
numerically ineffective at accessible scales. Thus this note
isolates a strong obstruction; it is \emph{not} yet a breakthrough
and should not be integrated in the main article as one.
\end{remark}
\end{document}
